% Generated draft: requires mathematical and source review.
\documentclass[11pt]{article}
\usepackage[margin=25mm]{geometry}
\usepackage{fontspec}
\setmainfont{Latin Modern Roman}
\usepackage{amsmath,amssymb,mathtools,unicode-math}
\setmathfont{Latin Modern Math}
\usepackage{enumitem,needspace,xurl,hyperref}
\setlength{\parindent}{0pt}
\setlength{\parskip}{5pt}
\setlength{\emergencystretch}{3em}
\newcommand{\N}{\mathbb{N}}
\newcommand{\Z}{\mathbb{Z}}
\newcommand{\Q}{\mathbb{Q}}
\newcommand{\R}{\mathbb{R}}
\newcommand{\C}{\mathbb{C}}
\newcommand{\F}{\mathbb{F}}
\newcommand{\A}{\mathbb{A}}
\newcommand{\PP}{\mathbb{P}}
\newcommand{\E}{\mathbb{E}}
\newcommand{\eps}{\varepsilon}
\newcommand{\dd}{\,\mathrm{d}}
\begin{document}
\section*{UnsolvedMath ID: 2308004}
\textbf{Source problem code:} catalogue entry
\subsection{Definitions and mathematical statement}
For a nonempty connected open set $D\subseteq\mathbb C$, let
\[
A(D)=\{f:D\to\mathbb C:f\text{ is holomorphic}\}
\]
with pointwise addition and multiplication. Initially $D$ is simply connected; $A(D)$ does not impose continuity on $\partial D$.

\textbf{Questions (J. J. Kelleher).} Can $A(D)$ be the direct sum of two nontrivial subrings of itself? Is there a more general result for multiply connected domains $D$? The source describes the proposed result as a generalization of Taylor's theorem but does not specify which theorem or decomposition is meant.

There are distinct algebraic readings of direct sum. A ring direct-sum decomposition $A(D)=B\oplus C$ requires unique additive representations and componentwise multiplication, in particular $BC=CB=\{0\}$. An additive direct sum of subrings requires only $A(D)=B+C$ and $B\cap C=\{0\}$, allowing nonzero cross products. Requiring both subrings to contain the same identity is another restriction. The source does not choose among these readings or define nontrivial (for example, whether constant subrings or evaluation ideals are excluded). No topological closure or complex-linearity condition is supplied. These choices are essential and are not silently fixed.

The source also attributes to Bers, expressly without a citation, that the algebraic structure of $A(D)$ determines the conformal structure of $D$, even for arbitrary connectivity. This is contextual attribution; the type of ring or algebra isomorphism intended is not specified. The supplied summary reports the decomposition question open as of the 2018 edition, with no definitive resolution located.
\subsection{Short English statement}
Ask whether a holomorphic-function ring can split into two nontrivial subrings, and seek an extension to multiply connected domains. The kind of direct sum, allowed subrings, and referenced Taylor theorem are unspecified.
\subsection{Sources}
\begin{itemize}
\item Hayman, \emph{Research Problems in Function Theory} (2018), Problem 8.4, attributed to J. J. Kelleher, \url{https://arxiv.org/abs/1809.07200}; source-TeX extraction, accessed 2026-07-29. The 2018 update was retained for status triage.
\item Bers, the stated determination of conformal structure, explicitly with no citation in the record; Taylor's theorem, named without a formulation or reference. No bibliographic details are supplied for either attribution.
\item Local registry \texttt{data/unsolvedmath/problems.json}, permanent ID $2308004$, source code catalogue entry, \emph{Research Problems in Function Theory --- Problem 8.4}. The record labels this open and OPEN-TRIAGE, with NEEDS\_REVIEW source-status and rights qualifications.
\end{itemize}
\end{document}
